\documentclass[11pt]{article}
\usepackage{amsmath,amssymb,amsthm,graphicx,booktabs,hyperref}
\newtheorem{theorem}{Theorem}
\newtheorem{lemma}{Lemma}
\newcommand{\norm}[1]{\left\lVert #1 \right\rVert}
\newcommand{\E}{\mathbb{E}}
\title{Blockwise Sparse Attention with Provable Error Bounds}
\author{Ada Smith\thanks{Corresponding author.} \and Bo Chen \and Chidi Okafor}
\date{November 2026}
\begin{document}
\maketitle
\begin{abstract}
We analyse a blockwise sparsity pattern for self-attention, prove an error bound against dense attention, and measure the speed-up on long inputs.
\end{abstract}
\section{Introduction}
Self-attention costs $O(n^2)$ in the sequence length $n$. We restrict each query to a set of blocks and show that the output error is bounded by the discarded softmax mass. Our contributions are:
\begin{enumerate}
  \item a block pattern with $O(n\sqrt{n})$ cost;
  \item Theorem~\ref{thm:bound}, an error bound;
  \item an evaluation on three tasks (Table~\ref{tab:main}).
\end{enumerate}
\section{Setup}\label{sec:setup}
Let $Q,K,V\in\mathbb{R}^{n\times d}$. Dense attention is
\begin{equation}\label{eq:dense}
  A = \mathrm{softmax}\!\left(\frac{QK^\top}{\sqrt d}\right)V .
\end{equation}
\subsection{Blocks}
Partition $\{1,\dots,n\}$ into blocks of size $b$. A \emph{mask} $M\in\{0,1\}^{n\times n}$ keeps entry $(i,j)$ iff block$(j)$ is selected for block$(i)$.
\section{Analysis}
\begin{lemma}\label{lem:mass}
Let $p_i$ be the softmax row. If the selected blocks hold mass at least $1-\epsilon$ then $\norm{p_i - \tilde p_i}_1 \le 2\epsilon$.
\end{lemma}
\begin{proof}
The renormalised row $\tilde p_i$ differs from $p_i$ only by the discarded mass $\epsilon$ and the rescaling of the kept mass. Summing the absolute differences gives the claim.
\end{proof}
\begin{theorem}\label{thm:bound}
Under Lemma~\ref{lem:mass}, $\norm{A-\tilde A}_{\infty} \le 2\epsilon \max_j \norm{V_j}$.
\end{theorem}
\section{Experiments}
\begin{table}[h]
\centering
\caption{Accuracy (\%) and relative speed.}
\label{tab:main}
\begin{tabular}{lccc}
\toprule
Model & Task A & Task B & Speed \\
\midrule
Dense & 81.2 & 77.4 & 1.0$\times$ \\
Blockwise & 80.9 & 77.1 & 3.4$\times$ \\
\bottomrule
\end{tabular}
\end{table}
The results in Table~\ref{tab:main} follow Equation~\eqref{eq:dense}.\footnote{Code: \url{https://example.org/blocks}.}
\section*{Acknowledgements}
We thank the reviewers.
\end{document}
